%!TEX root = ../dokumentation.tex

%
% Nahezu alle Einstellungen koennen hier getaetigt werden
%
\RequirePackage{luatex85}
\RequirePackage[l2tabu, orthodox]{nag}    % weist in Commandozeile bzw. log auf veraltete LaTeX Syntax hin

\documentclass[%
    lualatex,
    oneside,            % Einseitiger Druck.
    12pt,                % Schriftgroesse
    parskip=half,        % Halbe Zeile Abstand zwischen Absätzen.
%	topmargin = 10pt,	% Abstand Seitenrand (Std:1in) zu Kopfzeile [laut log: unused]
    headheight = 12pt,    % Höhe der Kopfzeile
%	headsep = 30pt,	% Abstand zwischen Kopfzeile und Text Body  [laut log: unused]
    headsepline,        % Linie nach Kopfzeile.
    footsepline,        % Linie vor Fusszeile.
    footheight = 16pt,    % Höhe der Fusszeile
    abstracton,        % Abstract Überschriften
    DIV=calc,        % Satzspiegel berechnen
    BCOR=8mm,        % Bindekorrektur links: 8mm
    headinclude=false,    % Kopfzeile nicht in den Satzspiegel einbeziehen
    footinclude=false,    % Fußzeile nicht in den Satzspiegel einbeziehen
    listof=totoc,        % Abbildungs-/ Tabellenverzeichnis im Inhaltsverzeichnis darstellen
    toc=bibliography,    % Literaturverzeichnis im Inhaltsverzeichnis darstellen
]{scrreprt}    % Koma-Script report-Klasse, fuer laengere Bachelorarbeiten alternativ auch: scrbook

% Einstellungen laden
\usepackage{xstring}
\usepackage{fontspec}
\setmainfont{SWRSans-Regular.ttf}[
    Path=font/SWR_Sans/,
    BoldFont=SWRSans-SemiBold.ttf,
    ItalicFont=SWRSans-Italic.ttf,
    BoldItalicFont=SWRSans-SemiBoldItalic.ttf
]

\newcommand{\einstellung}[1]{%
    \expandafter\newcommand\csname #1\endcsname{}
    \expandafter\newcommand\csname setze#1\endcsname[1]{\expandafter\renewcommand\csname#1\endcsname{##1}}
}
\newcommand{\langstr}[1]{\einstellung{lang#1}}

\input{ads/einstellungen_liste} % verfügbare Einstellungen
\input{einstellungen} % lese Einstellungen

\input{lang/strings} % verfügbare Strings
\setzelangdeckblattabschlusshinleitung{für die Prüfung zum}
\setzelangartikelstudiengang{des}
\setzelangstudiengang{Studiengangs}
\setzelanganderdh{an der Dualen Hochschule Baden-Württemberg}
\setzelangvon{von}
\setzelangdbbearbeitungszeit{Bearbeitungszeitraum}
\setzelangdbmatriknr{Matrikelnummer}
\setzelangdbkurs{Kurs}
\setzelangdbfirma{Ausbildungsfirma}
\setzelangdbbetreuer{Betreuer des Dualen Partners}
\setzelangdbgutachter{Gutachter der Dualen Hochschule}
\setzelangsperrvermerk{Sperrvermerk}
\setzelangerklaerung{Erklärung}
\setzelangabkverz{Abkürzungsverzeichnis}
\setzelangglossar{Glossar}
\setzelanganhang{Anhang}
\setzelangabstract{Abstract}
\setzelanglistingname{Listing}
\setzelanglistlistingname{Listings}
\setzelanglistingautorefname{Listing} % Übersetzung einlesen

% Einstellung der Sprache des Paketes Babel und der Verzeichnisüberschriften
\usepackage[english, ngerman]{babel}

%%%%%%% Package Includes %%%%%%%
\usepackage[margin=\seitenrand,foot=1cm]{geometry}    % Seitenränder und Abstände
\usepackage[activate]{microtype} %Zeilenumbruch und mehr
\usepackage[onehalfspacing]{setspace}
\usepackage{makeidx}
\usepackage[autostyle=true,german=quotes]{csquotes}
\usepackage{longtable}
\usepackage{enumitem}    % mehr Optionen bei Aufzählungen
\usepackage{graphicx}
\usepackage{pdfpages}   % zum Einbinden von PDFs
\usepackage{xcolor}    % für HTML-Notation
\usepackage{float}
\usepackage{array}
\usepackage{calc}        % zum Rechnen (Bildtabelle in Deckblatt)
\usepackage[right]{eurosym}
\usepackage{wrapfig}
%\usepackage{pgffor} % für automatische Kapiteldateieinbindung
\usepackage[perpage, hang, multiple, stable]{footmisc} % Fussnoten
\usepackage[printonlyused, footnote]{acronym} % falls gewünscht kann die Option footnote eingefügt werden, dann wird die Erklärung nicht inline sondern in einer Fußnote dargestellt
\usepackage{listings}

% Notizen. Einsatz mit \todo{Notiz} oder \todo[inline]{Notiz}. 
\usepackage[obeyFinal,backgroundcolor=yellow,linecolor=black]{todonotes}
% Alle Notizen ausblenden mit der Option "final" in \documentclass[...] oder durch das auskommentieren folgender Zeile
% \usepackage[disable]{todonotes}

% Kommentarumgebung. Einsatz mit \comment{}. Alle Kommentare ausblenden mit dem Auskommentieren der folgenden und dem aktivieren der nächsten Zeile.
\newcommand{\comment}[1]{\par {\bfseries \color{blue} #1 \par}} %Kommentar anzeigen
% \newcommand{\comment}[1]{} %Kommentar ausblenden


%%%%%% Configuration %%%%%

%% Anwenden der Einstellungen

\ladefarben{}

% Titel, Autor und Datum
\title{\titel}
\author{\autor}
\date{\datum}

% PDF Einstellungen
\usepackage[%
    pdftitle={\titel},
    pdfauthor={\autor},
    pdfsubject={\arbeit},
    pdfcreator={LaTeX},
    pdfkeywords={\arbeit, \autor, \titel},
    pdfproducer={LaTeX},
    pdfpagemode=UseOutlines,         % Beim Oeffnen Inhaltsverzeichnis anzeigen
    pdfdisplaydoctitle=true,                % Dokumenttitel statt Dateiname anzeigen.
    pdftoolbar=false,                           % Acrobat-Toolbar anzeigen.
    pdfmenubar=false,                        % Acrobat-Menü anzeigen.
    pdflang={\sprache},                       % Sprache des Dokuments.
]{hyperref}

% (Farb-)einstellungen für die Links im PDF
\hypersetup{%
    colorlinks=true,        % Aktivieren von farbigen Links im Dokument
    linkcolor=black,    % Farbe festlegen
    citecolor=LinkColor,
    filecolor=LinkColor,
    menucolor=LinkColor,
    urlcolor=LinkColor,
    linktocpage=false,        % Nicht der Text sondern die Seitenzahlen in Verzeichnissen klickbar
    bookmarksnumbered=true    % Überschriftsnummerierung im PDF Inhalt anzeigen.
}
% Workaround um Fehler in Hyperref, muss hier stehen bleiben
\usepackage{bookmark} %nur ein latex-Durchlauf für die Aktualisierung von Verzeichnissen nötig

% Schriftart in Captions etwas kleiner
\addtokomafont{caption}{\small}

% Literaturverweise (sowohl deutsch als auch englisch)

\usepackage[
    backend=biber,        % empfohlen. Falls biber Probleme macht: bibtex
    bibwarn=true,
    bibencoding=utf8,    % wenn .bib in utf8, sonst ascii
    sortlocale=de_DE,
    style=\zitierstil,
]{biblatex}


\ladeliteratur{}

%%%%%% Additional settings %%%%%%

% Hurenkinder und Schusterjungen verhindern
% http://projekte.dante.de/DanteFAQ/Silbentrennung
\clubpenalty = 10000 % schließt Schusterjungen aus (Seitenumbruch nach der ersten Zeile eines neuen Absatzes)
\widowpenalty = 10000 % schließt Hurenkinder aus (die letzte Zeile eines Absatzes steht auf einer neuen Seite)
\displaywidowpenalty=10000

% Bildpfad
%\graphicspath{{images/}}

% Einige häufig verwendete Sprachen
\lstloadlanguages{PHP,Python,Java,C,C++,bash}
\listingsettings{}
% Umbennung des Listings
\renewcommand\lstlistingname{\langlistingname}
\renewcommand\lstlistlistingname{\langlistlistingname}
\newcommand{\lstlistingauacrfacrotorefname}{\langlistingautorefname}

% Abstände in Tabellen
\setlength{\tabcolsep}{\spaltenabstand}
\renewcommand{\arraystretch}{\zeilenabstand}

% Eigene Pakete
\usepackage{adjustbox}
\usepackage{svg}
\usepackage{xurl}
\usepackage{tikz}
\usepackage{tabularx}
\usepackage[table]{xcolor}